% =====================================================================
\chapter{La máscara causal}
\label{tema:3}
% =====================================================================

Un modelo de lenguaje autorregresivo factoriza la probabilidad de una secuencia con la regla de la cadena (módulo~1, tema~2): la distribución del token $t+1$ solo puede depender de los tokens $1,\dots,t$. Este tema define las máscaras aditivas, demuestra que sumar $-\infty$ antes de la softmax equivale a restringir la softmax al conjunto de posiciones permitidas, prueba que la máscara causal hace que cada salida dependa solo de su pasado y deduce de ello que una única pasada produce las $n$ predicciones de una secuencia. Se siguen la §3.2.3 de~\cite{vaswani}, la §7.3 de~\cite{jm3}, las §3.4.4 y §3.4.5 de la tarea~1 de CS336~\cite{cs336} y la §11.3.2.1 de~\cite{d2l}.

\section{Máscaras aditivas}

\begin{definicion}[Máscara aditiva]\label{def:mascara}
Sea $\mathcal{P}\subseteq\{1,\dots,n\}\times\{1,\dots,m\}$ el conjunto de pares $(i,j)$ \emph{permitidos}, con la condición de que cada fila $i$ contenga al menos un par permitido. La \emph{máscara aditiva} asociada es la matriz $M\in\{0,-\infty\}^{n\times m}$ con $M_{ij}=0$ si $(i,j)\in\mathcal{P}$ y $M_{ij}=-\infty$ en otro caso. La \emph{atención enmascarada} es
\[
\Att_M(Q,K,V)=\softmax\!\Bigl(\frac{QK\T}{\sqrt{d_k}}+M\Bigr)V,
\]
con el convenio $e^{-\infty}=0$.
\end{definicion}

\begin{definicion}[Máscara causal]\label{def:causal}
En la autoatención ($m=n$), la \emph{máscara causal} es la asociada a $\mathcal{P}=\{(i,j):j\le i\}$:
\[
M_{ij}=\begin{cases}0,& j\le i,\\-\infty,& j>i.\end{cases}
\]
Es triangular estrictamente superior en sus entradas $-\infty$; la diagonal siempre está permitida.
\end{definicion}

Vaswani \emph{et al.} la describen como el bloqueo de las conexiones «ilegales» de la autoatención del decodificador, fijando a $-\infty$ las entradas correspondientes de la softmax~\cite[§3.2.3]{vaswani}; Jurafsky y Martin la presentan en la §7.3 de~\cite{jm3}, y la tarea~1 de CS336 la implementa con una máscara booleana en la que \texttt{True} indica «puede atender»~\cite[§3.4.4]{cs336}. La \cref{fig:causal} muestra su estructura.

\addtocounter{figure}{1}%

\begin{proposicion}[Softmax restringida]\label{prop:restringida}
Sea $z\in\R^{m}$ y sea $\mathcal{P}_i\subseteq\{1,\dots,m\}$ no vacío. Si $\mu_j=0$ para $j\in\mathcal{P}_i$ y $\mu_j=-\infty$ en otro caso, entonces
\[
\softmax(z+\mu)_j=\begin{cases}\dfrac{e^{z_j}}{\sum_{l\in\mathcal{P}_i}e^{z_l}},& j\in\mathcal{P}_i,\\[2ex] 0,& j\notin\mathcal{P}_i.\end{cases}
\]
En particular, los pesos permitidos suman~1 y no dependen de las puntuaciones $z_j$ con $j\notin\mathcal{P}_i$.
\end{proposicion}

La condición $\mathcal{P}_i\ne\varnothing$ es necesaria: una fila enteramente enmascarada daría $0/0$, que en coma flotante produce \texttt{NaN} (módulo~11). En la máscara causal no puede ocurrir, porque $(i,i)$ siempre está permitido.

\begin{observacion}[Por qué $-\infty$ antes de la softmax]
Sumar~0 antes de la softmax no anula nada, porque $e^{0}=1$. Poner a~0 los pesos \emph{después} de la softmax deja filas que suman menos de~1 y pesos del pasado más pequeños de lo debido; hay que renormalizar, y el resultado coincide entonces con el de la \cref{prop:restringida} (\cref{eje:3-renorm}). En la práctica se usa \texttt{-inf} o un número negativo grande; el \cref{eje:3-finito} cuantifica cuánto debe valer.
\end{observacion}

\addtocounter{definicion}{1}%

\section{Causalidad}

\begin{teorema}[Causalidad de la atención enmascarada]\label{teo:causal}
Sean $q_i=x_i\WQ$, $k_j=x_j\WK$, $v_j=x_j\WV$ y sea $o_i$ la fila $i$ de la autoatención con máscara causal. Entonces, para cada $i$, $o_i$ es función únicamente de $x_1,\dots,x_i$.
\end{teorema}

\begin{corolario}[Entrenamiento en paralelo]\label{cor:paralelo}
Sea $f$ la autoatención con máscara causal y, para cada $t\le n$, sea $f_t$ la misma operación aplicada solo al prefijo $x_{1:t}$. Entonces la fila $t$ de $f(x_{1:n})$ coincide con la última fila de $f_t(x_{1:t})$. En consecuencia, una sola evaluación sobre la secuencia completa proporciona las $n$ representaciones que una red usa para predecir $x_2,\dots,x_{n+1}$, cada una calculada solo con su prefijo.
\end{corolario}

La tarea~1 de CS336 describe la alternativa ingenua, ejecutar la atención $n$ veces, una por prefijo, y la descarta en favor de la máscara~\cite[§3.4.5]{cs336}. Una red recurrente necesita, en cambio, $n$ pasos secuenciales, porque el estado en $t$ depende del estado en $t-1$ (módulo~12). Esta diferencia, que Vaswani \emph{et al.} cuantifican como $O(1)$ operaciones secuenciales frente a $O(n)$~\cite[§4, tabla~1]{vaswani}, es la razón práctica principal por la que los Transformer desplazaron a las redes recurrentes en el preentrenamiento de modelos de lenguaje.

\begin{observacion}[Otras máscaras]
La \cref{def:mascara} admite otros conjuntos $\mathcal{P}$. La máscara de relleno (\emph{padding}) prohíbe las posiciones añadidas para igualar la longitud de las secuencias de un lote~\cite[§11.3.2.1]{d2l}; la de ventana deslizante permite solo $i-w<j\le i$, lo que reduce el coste a $O(nw)$ (\cref{tema:4}). La combinación de varias máscaras es la intersección de sus conjuntos permitidos, es decir, la suma de sus matrices.
\end{observacion}

